% ======================================================================
\section{Discussion and Limitations}
\label{sec:discussion}
% ======================================================================
\textbf{Scope of the privacy guarantees.} HALX guarantees
$k$-anonymity of the quasi-identifiers it stores, logs or forwards, and
$(\varepsilon,\delta)$-DP of the released normalization statistics. It does
\emph{not} make the encoder or the trees differentially private: a
white-box adversary could still attempt membership inference on the
weights. DP-SGD~\cite{abadi2016dpsgd} or federated training across
institutions~\cite{mcmahan2017fedavg,ruzafa2023fl} would close this gap.
$k$-anonymity also does not prevent attribute disclosure within an
equivalence class~\cite{sweeney2002kanon}.

\textbf{Benchmarks vs.\ bank traffic.} Neither benchmark is financial
traffic, and CICIDS2017 models transfer poorly to other
networks~\cite{catillo2022transferability}. Our CICIDS2017 subset (2\% of
the flows) is near the ceiling for all tree-based models, so the shifted
UNSW-NB15 split is more informative. Validation on bank flow logs is
planned.

\textbf{Temporal context.} HALX scores flows independently. Multi-stage
APT campaigns need cross-flow context (host-level session graphs or
windows over timestamped flows), which the public partitions used here,
lacking timestamps, cannot support without label leakage.

\textbf{Adversarial robustness.} We did not evaluate evasion attacks, which
\cite{arp2022dos} lists among the evaluation gaps in ML-based security. We
also did not test whether an attacker could manipulate the attention-based
explanations; both are left to future work.

% ======================================================================
\section{Conclusion}
\label{sec:conclusion}
% ======================================================================
We presented HALX, a hybrid Attention-LSTM and XGBoost detector for
financial network flows. Its privacy layer provides formally stated
$k$-anonymity and differential-privacy guarantees. In a controlled
comparison over three seeds, HALX matched the accuracy of the strongest
tree ensembles and clearly outperformed CNN-LSTM, BiLSTM and Transformer
detectors on CICIDS2017. Under the distribution shift of the official
UNSW-NB15 split it detected the most attacks at a 1\% false-alarm budget
(seed-0 ROC analysis), and it had the most stable false-positive rate
across seeds. The privacy layer
cost at most \PRIVCOSTMAX{} F1 points. The broader
lesson for financial SOCs is that boosted trees should make the final
decision on flow features. A deep encoder does not raise accuracy beyond
a tuned XGBoost, but it adds per-alert explanations, a more stable false-alarm
rate and a modest gain at low false-alarm budgets. Future work will add DP-SGD training,
federated deployment across institutions, adversarial evaluation and
validation on bank traffic.

